\section{Resumo}

Este relatório descreve \emph{como a tensão de um surto se distribui ao longo do enrolamento de um
gerador síncrono} com o neutro solidamente aterrado. Quando uma frente de onda rápida — o impulso
atmosférico padrão \SI{1.2}{\micro\second}/\SI{50}{\micro\second} — atinge a entrada da máquina, o
enrolamento não responde como uma única indutância concentrada: no primeiro instante ele se comporta
como uma \emph{rede capacitiva}, e a tensão se concentra de forma acentuada nas primeiras espiras.

O enrolamento é representado por uma escada de \Nsec{} seções com resistência e indutância série
(\Rsec, \Lsec), capacitância para o núcleo aterrado (\Cg, \emph{shunt}) e capacitância entre espiras
(\Cs, \emph{série}). A razão entre as duas capacitâncias define o parâmetro de distribuição
\(\adist = \sqrt{\Cg/\Cs}\), que controla quão desigual é a repartição inicial da tensão. Para a
máquina estudada adota-se \(\adist = 5\).

A história física tem dois instantes-chave. Em \(t = 0^+\), a distribuição é capacitiva e fortemente
não uniforme, \(v(x)/V_0 = \sinh\!\big(\adist(1-x)\big)/\sinh(\adist)\), com a entrada solicitada por
um fator \(\adist\coth\adist \approx 5\) acima da média. Depois, a indutância entra em ação e a
distribuição relaxa para um perfil quase uniforme, preso a zero no neutro aterrado. É essa
\emph{transição do perfil concentrado para o uniforme} que dimensiona o isolamento de entrada e que
este relatório documenta, com animações geradas a partir da simulação. Os resultados foram validados
contra a solução analítica clássica de surtos em enrolamentos \cite{greenwood1991} e por uma simulação
independente em ATP/EMTP \cite{dommel1969}.
